\documentclass[10pt,a4paper]{amsart}
\usepackage{amsmath,amssymb,amsthm,graphicx}
% Retained preamble for the Stacks Project chapter snapshot in this directory.
% The chapter's own preamble.tex is a copy of the shared Stacks Project preamble; the block
% below is sliced MECHANICALLY, line for line, from it. Nothing was added, removed or
% redefined; the wrapper class, the xr-hyper/hyperref wiring and the externaldocument list
% are omitted because the assembler supplies the document class.
%
% lines taken from preamble.tex: 3, 5, 9, 11, 12, 156, 157, 158, ... (51 lines total)
%
% package coverage restored from the chapter's own preamble (2026-10-13): xy, xy-2cell

\usepackage[all]{xy}
\xyoption{2cell}
\UseAllTwocells

\newenvironment{reference}{\comment}{\endcomment}
\newenvironment{slogan}{\comment}{\endcomment}
\newenvironment{history}{\comment}{\endcomment}
\theoremstyle{plain}
\newtheorem{theorem}[subsection]{Theorem}
\newtheorem{proposition}[subsection]{Proposition}
\newtheorem{lemma}[subsection]{Lemma}
\theoremstyle{definition}
\newtheorem{definition}[subsection]{Definition}
\newtheorem{example}[subsection]{Example}
\newtheorem{exercise}[subsection]{Exercise}
\newtheorem{situation}[subsection]{Situation}
\theoremstyle{remark}
\newtheorem{remark}[subsection]{Remark}
\newtheorem{remarks}[subsection]{Remarks}
\numberwithin{equation}{subsection}
\def\lim{\mathop{\mathrm{lim}}\nolimits}
\def\colim{\mathop{\mathrm{colim}}\nolimits}
\def\Spec{\mathop{\mathrm{Spec}}}
\def\Hom{\mathop{\mathrm{Hom}}\nolimits}
\def\Ext{\mathop{\mathrm{Ext}}\nolimits}
\def\SheafHom{\mathop{\mathcal{H}\!\mathit{om}}\nolimits}
\def\SheafExt{\mathop{\mathcal{E}\!\mathit{xt}}\nolimits}
\def\Sch{\mathit{Sch}}
\def\Mor{\mathop{\mathrm{Mor}}\nolimits}
\def\Ob{\mathop{\mathrm{Ob}}\nolimits}
\def\Sh{\mathop{\mathit{Sh}}\nolimits}
\def\NL{\mathop{N\!L}\nolimits}
\def\CH{\mathop{\mathrm{CH}}\nolimits}
\def\proetale{{pro\text{-}\acute{e}tale}}
\def\etale{{\acute{e}tale}}
\def\QCoh{\mathit{QCoh}}
\def\Ker{\mathop{\mathrm{Ker}}}
\def\Im{\mathop{\mathrm{Im}}}
\def\Coker{\mathop{\mathrm{Coker}}}
\def\Coim{\mathop{\mathrm{Coim}}}
\DeclareMathSymbol{\boxtimes}{\mathbin}{AMSa}{"02}
\def\QCohstack{\mathcal{QC}\!\mathit{oh}}
\def\Cohstack{\mathcal{C}\!\mathit{oh}}
\def\Spacesstack{\mathcal{S}\!\mathit{paces}}
\def\Quotfunctor{\mathrm{Quot}}
\def\Hilbfunctor{\mathrm{Hilb}}
\def\Curvesstack{\mathcal{C}\!\mathit{urves}}
\def\Polarizedstack{\mathcal{P}\!\mathit{olarized}}
\def\Complexesstack{\mathcal{C}\!\mathit{omplexes}}
\def\Pic{\mathop{\mathrm{Pic}}\nolimits}
\def\Picardstack{\mathcal{P}\!\mathit{ic}}
\def\Picardfunctor{\mathrm{Pic}}
\def\Deformationcategory{\mathcal{D}\!\mathit{ef}}
\usepackage{cleveref}
\crefname{originaldefinition}{original definition}{original definitions}
\Crefname{originaldefinition}{Original definition}{Original definitions}
\makeatletter
\newlabel{definition-algebraic-space}{{025Y}{source}{}{sourcepointer-1}{}}
\newlabel{definition-algebraic-space@cref}{{[originaldefinition][0][]025Y}{[1][0][]source}{}{}{}}
\makeatother
\crefname{originallemma}{original lemma}{original lemmas}
\Crefname{originallemma}{Original lemma}{Original lemmas}
\makeatletter
\newlabel{lemma-base-change-representable-transformations}{{02WB}{source}{}{sourcepointer-2}{}}
\newlabel{lemma-base-change-representable-transformations@cref}{{[originallemma][0][]02WB}{[1][0][]source}{}{}{}}
\makeatother
\crefname{originallemma}{original lemma}{original lemmas}
\Crefname{originallemma}{Original lemma}{Original lemmas}
\makeatletter
\newlabel{lemma-composition-representable-transformations}{{02WA}{source}{}{sourcepointer-3}{}}
\newlabel{lemma-composition-representable-transformations@cref}{{[originallemma][0][]02WA}{[1][0][]source}{}{}{}}
\makeatother
\crefname{originallemma}{original lemma}{original lemmas}
\Crefname{originallemma}{Original lemma}{Original lemmas}
\makeatletter
\newlabel{lemma-product-representable-transformations}{{02WC}{source}{}{sourcepointer-4}{}}
\newlabel{lemma-product-representable-transformations@cref}{{[originallemma][0][]02WC}{[1][0][]source}{}{}{}}
\makeatother
\title{A fibre product of algebraic spaces over a sheaf with representable diagonal}
\author{The Stacks Project Authors}
\date{stacks-project:spaces.tex:local-snapshot-20261012; GNU FDL 1.2 or later}
\begin{document}
\maketitle
\setcounter{subsection}{1}
\noindent\textit{Editorial note.} All delivered mathematical text is the author's own, copied from the pinned Stacks Project chapter file spaces.tex; no statement or proof step was rewritten, reordered or removed.\par\medskip
\noindent\textit{Editorial note.} The proof contains one xymatrix display; it is part of the counted proof span and is preserved verbatim, with the chapter's own Xy-pic setup retained in the preamble.\par\medskip
\noindent\textit{Editorial note.} The four references of the counted spans are recorded as non-load-bearing navigation pointers; the argument that consumes them stays inside the counted proof span.\par\medskip
\noindent\textit{Editorial note.} Preamble: sources/stacks-spaces/retained-preamble.tex, sliced mechanically from the chapter's own preamble.tex.\par\medskip
\noindent\textit{Allowed standard prerequisites (editorial).}\begin{itemize}
\item Definition of an algebraic space. Definition definition-algebraic-space, used at proof lines 1140-1141 and 1168-1170. Supplies the three conditions that are verified.
\item Products of representable transformations are representable. Lemma lemma-product-representable-transformations (409-420), used at proof lines 1152-1155. Makes Delta\_F x Delta\_G representable.
\item Base change of a representable transformation is representable. Lemma lemma-base-change-representable-transformations, used at proof lines 1152-1155 and 1163-1167. Carries representability to the diagonal and to the chart.
\item Composition of representable transformations with a property. Lemma lemma-composition-representable-transformations, used at proof lines 1163-1167. Gives the surjective etale chart.
\end{itemize}\medskip
\ifdefined\hypertarget\hypertarget{sourcepointer-1}{}\fi
\noindent\textit{Source reference.} Original definition 025Y: focused-seeds/AGP/sources/stacks-snapshot-20261011/tags.txt, source lines 2831--2831. Navigation pointer into the Stacks Project; the standard fact it supplies is applied inside the counted proof span and the argument resting on it stays there.\par\medskip
\ifdefined\hypertarget\hypertarget{sourcepointer-2}{}\fi
\noindent\textit{Source reference.} Original lemma 02WB: focused-seeds/AGP/sources/stacks-snapshot-20261011/tags.txt, source lines 3798--3798. Navigation pointer into the Stacks Project; the standard fact it supplies is applied inside the counted proof span and the argument resting on it stays there.\par\medskip
\ifdefined\hypertarget\hypertarget{sourcepointer-3}{}\fi
\noindent\textit{Source reference.} Original lemma 02WA: focused-seeds/AGP/sources/stacks-snapshot-20261011/tags.txt, source lines 3797--3797. Navigation pointer into the Stacks Project; the standard fact it supplies is applied inside the counted proof span and the argument resting on it stays there.\par\medskip
\ifdefined\hypertarget\hypertarget{sourcepointer-4}{}\fi
\noindent\textit{Source reference.} Original lemma 02WC: focused-seeds/AGP/sources/stacks-snapshot-20261011/tags.txt, source lines 3799--3799. Navigation pointer into the Stacks Project; the standard fact it supplies is applied inside the counted proof span and the argument resting on it stays there.\par\medskip
\begin{lemma}
\label{lemma-fibre-product-spaces-over-sheaf-with-representable-diagonal}
Let $S$ be a scheme contained in $\Sch_{fppf}$.
Let $H$ be a sheaf on $(\Sch/S)_{fppf}$ whose diagonal
is representable. Let $F, G$ be algebraic spaces over $S$.
Let $F \to H$, $G \to H$ be maps of sheaves.
Then $F \times_H G$ is an algebraic space.
\end{lemma}

\begin{proof}
We check the 3 conditions of
Definition \ref{definition-algebraic-space}.
A fibre product of sheaves is a sheaf, hence $F \times_H G$ is a sheaf.
The diagonal of $F \times_H G$ is the left vertical arrow in
$$
\xymatrix{
F \times_H G \ar[r] \ar[d]_\Delta &
F \times G \ar[d]^{\Delta_F \times \Delta_G} \\
(F \times F) \times_{(H \times H)} (G \times G) \ar[r] &
(F \times F) \times (G \times G)
}
$$
which is cartesian. Hence $\Delta$ is representable as the base change
of the morphism on the right which is representable, see
Lemmas \ref{lemma-product-representable-transformations} and
\ref{lemma-base-change-representable-transformations}.
Finally, let $U, V \in \Ob((\Sch/S)_{fppf})$
and $a : U \to F$, $b : V \to G$ be surjective and \'etale.
As $\Delta_H$ is representable, we see that $U \times_H V$ is a scheme.
The morphism
$$
U \times_H V \longrightarrow F \times_H G
$$
is surjective and \'etale as a composition of the base changes
$U \times_H V \to U \times_H G$ and $U \times_H G \to F \times_H G$
of the \'etale surjective morphisms $U \to F$ and $V \to G$, see
Lemmas \ref{lemma-composition-representable-transformations} and
\ref{lemma-base-change-representable-transformations}.
This proves the last condition of
Definition \ref{definition-algebraic-space}
holds and we conclude that $F \times_H G$ is an algebraic space.
\end{proof}

\end{document}
